\subsection{COMPARISON RELATIONS: II. STRONG RELATIONS}

DEFINITION 3. Given two real functions \(f, g\) belonging to \(\mathcal{H}(\mathfrak{F}, \mathbf{R})\) , which are \(\geqslant 0\) on a set in \(\mathfrak{F}\) , one says that \(f\) is negligible relative to \(g\) , or that \(g\) is preponderant over \(f\) (along \(\mathfrak{F}\) ), and one writes \(f \ll g\) or \(g \gg f\) if, for every \(\varepsilon > 0\) , there exists a set \(X \in \mathfrak{F}\) such that \(f(t) \leqslant \varepsilon g(t)\) for every \(t \in X\) .

Given two normed spaces \(V_{1}\) , \(V_{2}\) and two functions \(f_{1}\) , \(f_{2}\) belonging to \(\mathcal{H}(\mathfrak{F}, V_{1})\) and \(\mathcal{H}(\mathfrak{F}, V_{2})\) respectively, one says that \(f_{1}\) is negligible relative to \(f_{2}\) (along \(F\) ), and one writes \(f_{1} \ll f_{2}\) or \(f_{2} \gg f_{1}\) , if \(\|f_{1}\| \ll \|f_{2}\|\) .

For every scalar \(\lambda \neq 0\) the relation \(\mathbf{f}_1 \ll \mathbf{f}_2\lambda\) is equivalent to \(\mathbf{f}_1 \ll \mathbf{f}_2\) . The relation \(\mathbf{f}_1 \ll \mathbf{f}_2\) implies \(\mathbf{f}_1 \preccurlyeq \mathbf{f}_2\) but is not equivalent to it.

Note that the relation \(f_{1} \preccurlyeq f_{2}\) by no means implies the relation `` \(f_{1} \ll f_{2}\) or \(f_{1} \asymp f_{2}\) '': one has \(\sin x \preccurlyeq 1\) as x tends to \(+\infty\) , but neither of the relations \(\sin x \ll 1\) or \(\sin x \asymp 1\) is true.

Examples. 1) The relation \(\mathbf{f} \ll 1\) means that \(\mathbf{f}\) tends to 0 along \(\mathfrak{F}\) .

\begin{enumerate}
\def\labelenumi{\arabic{enumi})}
\setcounter{enumi}{1}
\item
  When \(\alpha\) and \(\beta\) are two real numbers such that \(\alpha < \beta\) one has \(x^{\alpha} \ll x^{\beta}\) as \(x\) tends to \(+\infty\) . Similarly, when \(m\) and \(n\) are two rational integers such that \(m < n\) one has \(z^{m} \ll z^{n}\) as the complex number \(z\) tends to \(\infty\) .
\item
  As \(x\) tends to \(+\infty\) one has \(x^n \ll e^x\) for every integer \(n\) (III, p.~105).
\item
  In \(R^{2}\) one has, as \((x, y)\) tends to \((0, 0)\) ,
\end{enumerate}

\[
x ^ {2} + y ^ {2} \ll | x | + | y |.
\]

The following propositions can be deduced immediately from def. 3:

PROPOSITION 4. If \(f, g, h\) are three functions in \(\mathcal{H}(\mathfrak{F}, \mathbf{R})\) then the relations \(f \preccurlyeq g\) and \(g \ll h\) (resp. \(f \ll g\) and \(g \preccurlyeq h\) ) imply \(f \ll h\) .

PROPOSITION 5. Let \(\mathbf{f}_1, \mathbf{f}_2\) be two functions in \(\mathcal{H}(\mathfrak{F}, V)\) and \(g\) a function in \(\mathcal{H}(\mathfrak{F}, \mathbf{R})\) . Then the relations \(\mathbf{f}_1 \ll g\) and \(\mathbf{f}_2 \ll g\) imply \(\mathbf{f}_1 + \mathbf{f}_2 \ll g\) .

Moreover, the same argument as in prop. 3 of V, p.~213 shows that:

PROPOSITION 6. With the notation of prop. 3, the relations \(\mathbf{f}_1 \ll g_1\) and \(\mathbf{f}_2 \preccurlyeq g_2\) (resp. \(\mathbf{f}_1 \preccurlyeq g_1\) and \(\mathbf{f}_2 \ll g_2\) ) imply \([\mathbf{f}_1.\mathbf{f}_2] \ll g_1g_2\) .

Prop. 4 shows that the relation \(\mathbf{f} \ll \mathbf{g}\) between functions in \(\mathcal{H}(\mathfrak{F}, V)\) is transitive; but it is not reflexive: to be precise, the relation \(\mathbf{f} \ll \mathbf{f}\) implies that \(\mathbf{f}\) vanishes on a set in \(\mathfrak{F}\) (in other words, \(\mathbf{f}\) is equivalent to 0 modulo \(R_{\infty}\) ); indeed, for an \(\varepsilon\) such that \(0 < \varepsilon < 1\) there exists an \(X \in \mathfrak{F}\) such that \(\| \mathbf{f}(x) \| \leqslant \varepsilon \| \mathbf{f}(x) \|\) for every \(x \in X\) , which is possible only if \(\mathbf{f}(x) = 0\) for every \(x \in X\) . It follows that the relation `` \(\mathbf{f} \ll \mathbf{g}\) and \(\mathbf{g} \ll \mathbf{f}\) '' is transitive and symmetric, but not reflexive: so it is not an equivalence relation (it implies that there is a set \(X \in \mathfrak{F}\) such that \(\mathbf{f}(x) = \mathbf{g}(x) = 0\) for every \(x \in X\) ).

PROPOSITION 7. If \(\mathbf{f}\) and \(\mathbf{g}\) are two functions in \(\mathcal{H}(\mathfrak{F}, V)\) then the relation \(\mathbf{f} - \mathbf{g} \ll \mathbf{f}\) is equivalent to \(\mathbf{f} - \mathbf{g} \ll \mathbf{g}\) .

Indeed, \(\mathbf{f} - \mathbf{g} \ll \mathbf{f}\) implies that for every \(\varepsilon > 0\) there exists an \(X \in \mathfrak{F}\) such that \(\| \mathbf{f}(x) - \mathbf{g}(x) \| \leqslant \varepsilon \| \mathbf{f}(x) \|\) for every \(x \in X\) . But then \((1 - \varepsilon) \| \mathbf{f}(x) \| \leqslant \| \mathbf{g}(x) \|\) , and consequently \(\mathbf{f} \preccurlyeq \mathbf{g}\) , whence (V, p.~215, prop. 4) \(\mathbf{f} - \mathbf{g} \ll \mathbf{g}\) .

COROLLARY. The relation \(f - g \ll f\) is an equivalence relation on \(\mathcal{H}(\mathfrak{F}, V)\) .

Indeed, if \(f - g \ll f\) and \(g - h \ll g\) then \(f - g \ll g\) , whence (V, p.~215, prop. 5) \(f - h \ll g\) , and consequently, since \(g \preccurlyeq f\) , we have \(f - h \ll f\) , which shows that this relation is transitive; it is symmetric by prop. 7, and is clearly reflexive, whence the corollary.

DEFINITION 4. Given two functions f, g of \(\mathcal{H}(\mathfrak{F}, V)\) one says that f and g are equivalent (along F), and writes \(f \sim g\) , if \(f - g \ll f\) .

The relation \(\mathbf{f} \sim \mathbf{g}\) implies \(\mathbf{f} \asymp \mathbf{g}\) but is not equivalent to it.

Examples. 1) If \(\mathbf{a}\) is a constant function, \(\neq 0\) on \(\mathbf{E}\) , the relation \(\mathbf{f} \sim \mathbf{a}\) means that \(\mathbf{f}\) tends to \(\mathbf{a}\) along \(\mathfrak{F}\) .

\begin{enumerate}
\def\labelenumi{\arabic{enumi})}
\setcounter{enumi}{1}
\item
  Let V be a normed space over a non-discrete valued field K, and let \(\mathbf{f}(x)=\mathbf{a}_{0}x^{n}+\mathbf{a}_{1}x^{n-1}+\cdots+\mathbf{a}_{n}\) be a polynomial in the variable \(x\in K\) , with coefficients in V, such that \(a_{0}\neq0\) . Then \(f(x)\sim a_{0}x^{n}\) as \(|x|\) tends to \(+\infty\) .
\item
  When the real number \(x\) tends to \(+\infty\) one has \(\left(1 + \frac{1}{x}\right)\sin x \sim \sin x\) .
\item
  As the complex variable \(z\) tends to 0 one has \(e^z - 1 \sim z\) . More generally, if \(V\) is a normed space over a valued field \(K\) , and \(\mathbf{f}\) is a function defined on a neighbourhood of \(x_0 \in K\) , with values in \(V\) , and admitting a derivative \(\mathbf{f}'(x_0) \neq 0\) at the point \(x_0\) , then, as \(x\) approaches \(x_0\) , we have \(\mathbf{f}(x) - \mathbf{f}(x_0) \sim \mathbf{f}'(x_0)(x - x_0)\) (I, p.~3, def. 1).
\item
  As \((x, y)\) approaches \((0, 0)\) in \(R^{2}\) one has
\end{enumerate}

\[
\sqrt {\sin^ {2} x + \sin^ {2} y} \sim \sqrt {x ^ {2} + y ^ {2}}.
\]

\begin{enumerate}
\def\labelenumi{\arabic{enumi})}
\setcounter{enumi}{5}
\item
  Let \(f(x, y)\) be a polynomial with real coefficients in two real variables x, y, not having a constant term. If, as x approaches 0 while remaining \textgreater0, there is a function \(\varphi(x)\) , continuous on an interval \([0, a]\) , and such that \(\varphi(0) = 0\) and \(f(x, \varphi(x)) = 0\) for \(0 \leqslant x \leqslant a\) , one can show that there are a rational number r and a real number \(\lambda \neq 0\) such that \(\varphi(x) \sim \lambda x^{r}\) (V, p.~259, exerc. 3).
\item
  For every \(x > 0\) let \(\pi(x)\) denote the number of prime numbers which are \(\leqslant x\) ; it has been proved that \(\pi(x) \sim x / \log x\) as \(x\) tends to \(+\infty\) .
\end{enumerate}

\pandocbounded{\includegraphics[keepaspectratio]{images/part002_abf04123.jpg}}

Remark. Note that the relation \(\mathbf{f} \sim \mathbf{g}\) by no means implies that the difference \(\mathbf{f} - \mathbf{g}\) tends to 0 along \(\mathfrak{F}\) ; this difference can even be unbounded, as is shown by the example \(x^{2} + x \sim x^{2}\) as \(x\) tends to \(+\infty\) .

PROPOSITION 8. Let K be a valued field and \(f_{1}, f_{2}, g_{1}, g_{2}\) four functions in \(\mathcal{H}(\mathfrak{F}, \mathrm{K})\) ; then the relations \(f_{1} \sim g_{1}\) and \(f_{2} \sim g_{2}\) imply \(f_{1}f_{2} \sim g_{1}g_{2}\) .

Indeed, we have \(f_{1}f_{2}-g_{1}g_{2}=f_{1}(f_{2}-g_{2})+(f_{1}-g_{1})g_{2}\) ; since \(f_{1}\preccurlyeq g_{1}\) , \(f_{1}-g_{1}\ll g_{1}\) and \(f_{2}-g_{2}\ll g_{2}\) , we have \(f_{1}f_{2}-g_{1}g_{2}\ll g_{1}g_{2}\) (V, p.~215, prop. 5 and 6).

In contrast, we have given an example in V, p.~214 where one has \(f_{1} = g_{1}\) , \(f_{2} \sim g_{2}\) and yet the relation \(f_{1} + f_{2} \asymp g_{1} + g_{2}\) does not hold (so neither a fortiori does

\[
f _ {1} + f _ {2} \sim g _ {1} + g _ {2}).
\]

The comparison relations \(f \ll g\) , \(f \sim g\) are called strong relations. Two functions f, g from \(\mathcal{H}(\mathfrak{F}, V)\) are called comparable (or strongly comparable when one wants to avoid any possible confusion) if they satisfy one of the three relations: \(f \ll g\) , \(f \gg g\) , or ``there exists a \(\lambda \neq 0\) such that \(f \sim g\lambda\) ''.

Remarks. 1) Two functions can be weakly comparable though not strongly comparable, for example the functions 1 and \(\sin x\) as \(x\) tends to \(+\infty\) .

\begin{enumerate}
\def\labelenumi{\arabic{enumi})}
\setcounter{enumi}{1}
\tightlist
\item
  In the definitions of the comparison relations \(f_{1} \preccurlyeq f_{2}\) and \(f_{1} \ll f_{2}\) the norms on the spaces \(V_{1}, V_{2}\) where \(f_{1}\) and \(f_{2}\) respectively take their values, are involved only apparently; in reality only the topologies on \(V_{1}\) and \(V_{2}\) are involved, for the relations \(f_{1} \preccurlyeq f_{2}\) and \(f_{1} \ll f_{2}\) are replaced by equivalent relations when one replaces the norm on \(V_{1}\) or \(V_{2}\) by an equivalent norm (Gen.~Top., IX, p.~170, def. 7).
\end{enumerate}

